\chapter{Desarrollo matemático}

\section{Algunos resultados a tener en cuenta}

En los capítulos 8, 9, y 10 del libro de texto~\cite{Wegner2024}
se presentan algunas características no directamente esperables de las
distribuciones más importantes en Estadística cuando se trabaja en espacios de alta dimensión.
Sin embargo, algunos fenómenos, como la concentración de medida o la ortogonalidad de vectores, pueden resultar
de utilidad en Análisis de Datos.

\section{Proyecciones aleatorias}

En esta sección introducimos el concepto de proyecciones aleatorias.
Sea $\left(\Omega,\Sigma,\mathbb{P}\right)$ un espacio de probabilidad,
en esta sección se introducen matrices aleatorias del tipo
\begin{equation*}
	U\colon\Omega\to\mathbb{R}^{k\times d},\quad
	U=
	\begin{bmatrix}
		U_{11} & \cdots & U_{1d} \\
		\vdots & \ddots & \vdots \\
		U_{k1} & \cdots & U_{kd}
	\end{bmatrix},
\end{equation*}
con $U_{ij}\sim\mathcal{N}\left(0,1\right)$, $i=1,\dotsc, k$ y $j=1,\dotsc,d$,
variables aleatorias normales unidimensionales independientes.
Esta matriz aleatoria será denotada por $U\sim\mathcal{N}\left(\symbf{0}_{k\times d},1\right)$.
Algunas propiedades que tiene la matriz aleatoria están basadas en el comportamiento de sus filas como vectores normales multivariantes.
Si se denotan las filas de $U$ como $U_{i}=\left[U_{i1},\dotsc,U_{id}\right]^{T}$, para $i=1,\dotsc,k$,
entonces todas ellas son vectores aleatorios independientes con distribución $\mathcal{N}\left(\symbf{0}_{d},I_{d}\right)$
(ver Proposición A.8) y se pueden explotar los teoremas  8.2, 8.3 y 8.4 del Capítulo 8 del libro de texto y,
a su vez, los resultados del Capítulo 10, cuando $d$ sea elevado.
Por ejemplo, si usamos el Teorema 10.7 de Ortogonalidad Gaussiana, se tiene que las filas de la matriz $U$ normalizadas son ``aproximadamente''
ortogonales unas a otras.
Entonces, normalizando la matriz aleatoria $U$, se obtendrá una matriz aleatoria $V$ tal que $VV^{T}≈I_{k}$.
Dicha propiedad nos recuerda a la propiedad de ortogonalidad en las matrices de proyección basadas en Álgebra lineal.
Una de las aplicaciones más conocidas de las proyecciones es que permiten transformar conjuntos de datos en $\mathbb{R}^{d}$
en el subespacio de $\mathbb{R}^{k}$ generado por las $k$ filas de la matriz ortogonal de proyección.
Se debe tener en cuenta que si $k<d$ entonces la proyección es una transformación lineal, que
cuenta con esta buena propiedad de ortonormalidad que evita la ``duplicidad'' de información en los datos proyectados,
pero que no necesariamente conserva las distancias entre puntos del espacio original $\mathbb{R}^{d}$.
Esta conservación de las distancias originales es una propiedad útil en contextos de clasificación,
por tanto supervisada como no supervisada, por ejemplo, en métodos como los $m$-vecinos más próximos.
El objetivo de este capítulo es probar que las proyecciones basadas en matrices aleatorias del tipo
$U\sim\mathcal{N}\left(\symbf{0}_{k\times d},1\right)$, además que ella tiene la característica de la
ortonormalidad, de forma aproximada, se mantienen también, de forma
aproximada, las distancias entre datos en el espacio original.

\begin{theorem}[Teorema de la Proyección aleatoria]
	Sea $\left(\Omega,\Sigma,\mathbb{P}\right)$ un espacio de
	probabilidad y, para $k<d$ sea $U\colon\Omega\to\mathbb{R}^{k\times d}$
	una matriz aleatoria con $U\sim\mathcal{N}\left(\symbf{0}_{k\times d},1\right)$.
	Para $\symbf{x}\in\mathbb{R}^{d}\setminus\left\{\symbf{0}_{d}\right\}$ y $0<\varepsilon\leq 1$ tenemos
	\begin{equation*}
		\mathbb{P}
		\left[
			\left|
			\left\|
			U\symbf{x}
			\right\|-
			\sqrt{k}
			\left\|
			\symbf{x}
			\right\|
			\right|\leq
			\varepsilon\sqrt{k}
			\left\|
			\symbf{x}
			\right\|
			\right]\geq
		1-2\exp\left(-k\frac{\varepsilon^{2}}{16}\right).
	\end{equation*}
\end{theorem}

\begin{proof}
	En primer lugar, se debe tener en cuenta que el producto matriz por
	vector de $U$ por $\mathbf{x} \in \mathbb{R}^d$ da como resultado un
	vector en $\mathbb{R}^k$ que contiene los productos escalares entre
	cada vector fila de la matriz $U$, denotados por $U_i$, con el vector $\mathbf{x}$. Es decir,
	$$
		U \mathbf{x}=\left[\left\langle U_1, \mathbf{x}\right\rangle, \ldots,\left\langle U_k, \mathbf{x}\right\rangle\right]^T .
	$$

	Ahora, fijado $\mathbf{x} \in \mathbb{R}^d \backslash\left\{\mathbf{0}_d\right\}$, se puede considerar el vector
	$U\frac{\mathbf{x}}{\|\mathbf{x}\|}$, cuya coordenada $i$-ésima es
	$$
		\left(U\frac{\mathbf{x}}{\|\mathbf{x}\|}\right)_i=\left\langle U_i, \frac{\mathbf{x}}{\|\mathbf{x}\|}\right\rangle=\sum_{j=1}^d U_{i j} \frac{x_i}{\|\mathbf{x}\|}
	$$
	Usando que $U_{i j} \sim \mathcal{N}(0,1)$ para todo $(i, j)$, se tiene que
	$\left(U\frac{\mathbf{x}}{\|\mathbf{x}\|}\right)_i \sim \mathcal{N}(0,1)$ (ver Proposición A. 4 luego, junto a la independencia de los $U_{i j}$, se tiene $U\frac{\mathbf{x}}{\|\mathbf{x}\|} \sim \mathcal{N}\left(\mathbf{0}_d, I_d\right)$ (ver Proposición A.8). Se puede utilizar ahora el Teorema 10.4 del libro de texto (Anillo Gaussiano) para obtener
	$$
		P[|\|U \mathbf{x}\|-\sqrt{k}\|\mathbf{x}\|| \leq \varepsilon \sqrt{k} \cdot\|\mathbf{x}\|]=P\left[\left|\left\|U \frac{\mathbf{x}}{\|\mathbf{x}\|}\right\|-\sqrt{k}\right| \leq \varepsilon \sqrt{k}\right] \geq 1-2\exp\left(k\frac{\varepsilon^{2}}{16}\right)
	$$
	En esta acotación se ha aplicado dicho teorema con $d=k$ y $\varepsilon=\varepsilon \sqrt{k}$.
	Nótese que se cumple la hipótesis $0<\varepsilon \sqrt{k} \leq \sqrt{k}$ dado que $0<\varepsilon \leq 1$.
\end{proof}
Puesto que $\sqrt{k}$ es fijo a lo largo de toda la prueba del Teorema 2.2.1, una reformulación más habitual de este resultado es
\begin{equation*}
	P[|\|\frac{1}{\sqrt{k}}U\frac{\mathbf{x}}{\|\mathbf{x}\|}\|-1| \leq \varepsilon]
	\geq 1-
	2\exp\left(-k\frac{\varepsilon^{2}}{16}\right)
\end{equation*}
\begin{definition}
	Sea $\left(\Omega,\Sigma,\mathbb{P}\right)$ un espacio de
	probabilidad y, para $k<d$, sea
	$U\colon\Omega\to\mathbb{R}^{k\times d}$ una matriz aleatoria con
	$U\sim\mathcal{N}\left(\symbf{0}_{k\times d},1\right)$.
	Además, sea $\omega\in\Omega$.
	Entonces, la aplicación
	\begin{align*}
		T_{U\left(\omega\right)}\colon\mathbb{R}^{d} & \longrightarrow
		\mathbb{R}^{k}                                                 \\
		\symbf{x}                                    & \longmapsto
		\frac{1}{\sqrt{k}}
		U\left(\omega\right)\symbf{x}
	\end{align*}
	se denomina proyección de Johnson-Lindenstrauss.
\end{definition}
Estrictamente, la transformación $T_{U\left(\omega\right)}$ no es una proyección en el sentido del
álgebra lineal, pues no necesariamente son matrices ortonormales.
Con esta notación, se puede reformular la conclusión del Teorema 2.3.1 como sigue:
\begin{equation*}
	\begin{split}
		\forall\symbf{x}\in\mathbb{R}^{d}\setminus\left\{\symbf{0}_{d}\right\}
		\text{ con }
		0<\varepsilon\leq 1
		\text{ y }
		U\sim\mathcal{N}\left(\symbf{0}_{k\times d},1\right): \\
		P[\left|\|T_{U\left(\omega\right)}\left(\mathbf{x}\right)\|-\|\mathbf{x}\|\right|\leq\varepsilon\|\mathbf{x}\|]
		\geq 1-
		2\exp\left(-k\frac{\varepsilon^{2}}{16}\right)
	\end{split},
\end{equation*}
se observa que con alta probabilidad
\begin{equation*}
	\frac{\left\|T_{U\left(\omega\right)}\left(\mathbf{x}\right)\right\|}{\left\|x\right\|}\approx 1.
\end{equation*}
Por tanto, de cierta forma, la proyección de Johnson-Lindenstrauss mantiene prácticamente invariante la norma de cada $\symbf{x}\in\mathbb{R}^{d}$.

\section{El teorema de Johnson-Lindenstrauss}

El siguiente teorema de Johnson-Lindenstrauss en el contexto probabilístico establece lo que sucede cuando consideramos
las distancias entre pares de puntos tras usar la proyección de Johnson-Lindenstrauss.

\begin{theorem}[Teorema de Johnson-Lindenstrauss]
	Sea $\left(\Omega,\Sigma,\mathbb{P}\right)$ un espacio de
	probabilidad y sea además $0<\varepsilon<1$, $n\geq 1$ y
	$k\geq\frac{48}{\varepsilon^2}\log n$.
	Sea $U\colon\Omega\to\mathbb{R}^{k\times d}$ una matriz aleatoria
	con $U\sim\mathcal{N}\left(\symbf{0}_{k\times d},1\right)$.
	Entonces, para cada
	$\left\{\symbf{x}_{1},\dotsc,\symbf{x}_{n}\right\}\subset\mathbb{R}^{d}$
	se cumple la siguiente estimación:
	\begin{equation*}
		\mathbb{P}
		\left[
			\forall i,j:
			\left(1-\varepsilon\right)
			\left\|
			\symbf{x}_{i}-
			\symbf{x}_{j}
			\right\|\leq
			\left\|
			T_{U\left(\omega\right)}\big(\symbf{x}_{i}\big)-
			T_{U\left(\omega\right)}\big(\symbf{x}_{j}\big)
			\right\|\leq
			\left(1+\epsilon\right)
			\left\|
			\symbf{x}_{i}-
			\symbf{x}_{j}
			\right\|
			\right]\geq
		1-\frac{1}{n}.
	\end{equation*}
\end{theorem}


\begin{proof}
	Para cada $i\neq j$ se define $\mathbf{x}=\mathbf{x}_i-\mathbf{x}_{j}$ y se aplica el Teorema 2.2.1 para acotar
	$\mathbb{P}_{i,j}$ como sigue:
	\begin{align*}
		\mathbb{P}_{i,j} & =
		\mathbb{P}
		\left[(1-\varepsilon)
			\left\|
			\mathbf{x}_{i}-
			\mathbf{x}_{j}
			\right\|\leq
			\left\|
			T_{U\left(\omega\right)}\left(\mathbf{x}_{i}\right)-
			T_{U\left(\omega\right)}\left(\mathbf{x}_{j}\right)\right\|
			\leq
			\left(1+\varepsilon\right)
			\left\|
			\mathbf{x}_{i}-
			\mathbf{x}_{j}
			\right\|
		\right]                                                                                                                                \\
		                 & =
		\mathbb{P}
		\left[
			\left(1-\varepsilon\right)
			\left\|\mathbf{x}\right\|\leq
			\left\|T_{U\left(\omega\right)}\left(\mathbf{x}\right)\right\|\leq
			\left(1+\varepsilon\right)
			\left\|\mathbf{x}\right\|
		\right]                                                                                                                                \\
		                 & =
		\mathbb{P}
		\left[
		\left|\left\|T_{U\left(\omega\right)}\left(\mathbf{x}\right)\right\|-\|\mathbf{x}\|\right| \leq \varepsilon \cdot\|\mathbf{x}\|\right] \\
		                 & \geq
		1-
		2\exp\left(-k\frac{\varepsilon^{2}}{16}\right).
		\tag{2.3}
	\end{align*}
	Puesto que hay $n$ puntos en el conjunto $\left\{\mathbf{x}_1, \ldots, \mathbf{x}_n\right\}$, existen $\binom{n}{2} \leq \frac{n^2}{2}$ pares distintos de puntos,
	por tanto, se puede acotar la probabilidad $\mathbb{P}$ teniendo en cuenta 2.3 como sigue:
	\begin{align*}
		 &
		\mathbb{P}
		\left[
			\left(1-\varepsilon\right)
			\left\|
			\mathbf{x}_{i}-
			\mathbf{x}_{j}
			\right\|\leq
			\left\|
			T_{U\left(\omega\right)}\left(\mathbf{x}_{i}\right)-
			T_{U\left(\omega\right)}\left(\mathbf{x}_{j}\right)
			\right\|\leq
			\left(1+\varepsilon\right)
			\left\|\mathbf{x}_{i}-\mathbf{x}_{j}\right\|
		\text { para todo } i,j\right]=                                                                                                                                                                                   \\
		 & =
		1-\mathbb{P}
		\left[\bigcup_{i, j}\left\{\left|\left\|T_U \mathbf{x}_i-T_U \mathbf{x}_{j}\right\|-\left\|\mathbf{x}_i-\mathbf{x}_{j}\right\|\right|>\varepsilon \cdot\left\|\mathbf{x}_i-\mathbf{x}_{j}\right\|\right\}\right]= \\
		 & =1-\sum_{i \neq j}\left(1-P_{i, j}\right) \geq 1-\sum_{i \neq j} 2 e^{-k \varepsilon^2 / 16} \geq 1-\frac{n^2}{2} 2 e^{-k \varepsilon^2 / 16} \geq                                                             \\
		 & \geq 1-n^2 \cdot e^{-\left(\frac{48}{\varepsilon^2} \log n\right) \varepsilon^2 / 16} \geq 1-n^2 \cdot e^{\log \left(n^{-3}\right)}=1-\frac{1}{n} \tag{5.4}
	\end{align*}
\end{proof}

\begin{remark}
	Para aplicar el Teorema 2.3.1 (de Johnson-Lindestrauss) correctamente se debe interpretar la relación
	que existe entre las distintas constantes que aparecen en el enunciado:
	$n$ tamaño de la muestra en consideración, $d$ dimensión original del conjunto de puntos y
	$k$ dimensión a la que se pretende proyectar los puntos.
	En particular, $k\geq\frac{48}{\varepsilon^{2}}\log n$ crece de forma logarítmica con respecto a $n$
	e inversamente proporcional al cuadrado de $\varepsilon$, lo cual quiere decir que conforme más grande sea el tamaño muestral $n$ y más pequeño
	sea $\varepsilon$ más grande será la dimensión $k$ a la que se debe proyectar.
\end{remark}

\begin{remark}
	La clave está en que la dimensión proyectada $k$ no depende de la dimensión original $d$, solo depende de $n$ y $\varepsilon$.
	El teorema de Johnson-Lindestrauss garantiza que, al proyectar un conjunto de $n$ puntos desde un espacio de
	alta dimensión a un subespacio euclídeo de dimensión $k\geq\frac{48}{\varepsilon^{2}}\log n$, es posible conservar las distancias entre todos los pares de puntos
	con un error de a lo sumo $\varepsilon$, y esto implica que la estructura geométrica del conjunto se
	preserva casi idénticamente por la proyección de Johnson-Lindestrauss $T_{U\left(\omega\right)}$.
\end{remark}

\begin{remark}
	El Teorema 2.3.1 proporciona resultados relevantes en el caso de que $d$ sea de una dimensión
	moderadamente grande en comparación con $k$ que es del orden de $\log n$.
	Esto nos permite incluso considerar conjuntos de puntos más grandes que la dimensión
	del espacio $d$ que se podrán proyectar a un espacio de dimensión menor sin perder su geometría.
\end{remark}

\begin{remark}
	Además, los resultados del Teorema 2.3.1 son totalmente independientes del conjunto original de puntos
	$\left\{\symbf{x}_{1},\dotsc,\symbf{x}_{n}\right\}\subset\mathbb{R}^{d}$
	o de la distribución que los generó (si la tuviesen).
	De hecho, únicamente son dependientes del tamaño del conjunto sobre el que se quieran aplicar.
\end{remark}

En este capítulo presentaremos el Teorema de Johnson-Lindenstrauss, que permitirá conservar
las buenas propiedades geométricas de los conjuntos de datos a la vez que se proyectan a un espacio de dimensión menor
donde las técnicas de Análisis de Datos clásicas se vuelven más fiables.

\chapter{Experimentos}
